\chapter{How to Keep Up}
\label{ch:keepingup}

You cannot. The field publishes faster than it can be read, models arrive
weekly, and tools turn over faster than their documentation. Accepting that is
more useful than any reading list: choose what to track, and accept gaps
everywhere else.

\section{What is stable}

Sort what you might learn by how long it stays true.

\begin{table}[htbp]
\centering
\caption{What to learn, and how long it lasts.}
\label{tab:halflife}
\small
\begin{tabularx}{\textwidth}{XlX}
\toprule
\textbf{Layer} & \textbf{Stays true for} & \textbf{Worth learning?} \\
\midrule
Model names and versions & Weeks & No. Look them up. \\
Benchmark scores & Months & No. They do not transfer to your task. \\
Prices and usage limits & Months & Only when buying. \\
Tool names and options & 6--12 months & Lightly; read the documentation. \\
Patterns (harness, RAG, agent loop) & Years & Yes. \\
Failures (hallucination, sycophancy, injection) & Years & Yes, thoroughly. \\
Data standards, provenance, reproducibility & Decades & Yes. This is what lasts. \\
Ecological expertise & A career & Yes. It is what makes you useful here. \\
\bottomrule
\end{tabularx}
\end{table}

\begin{keypoint}
Spend your learning time at the bottom of that table. Everything at the top is a
lookup.
\end{keypoint}

\section{People worth following}

A short list of people who are careful rather than fast.

\begin{description}
  \item[Melanie Mitchell (Santa Fe Institute)] The most reliable sceptic of AI
    capability claims \parencite{mitchell2026jagged,krakauer2026emergence}.
  \item[Blaise Ag\"uera y Arcas (Google)] Argues that prediction underlies both
    life and intelligence, so current models are intelligent in a real sense
    \parencite{aguerayarcas2025intelligence}. The strongest form of that case.
  \item[Anil Seth] A neuroscientist of consciousness who also puts prediction at
    the centre of the mind, but who is cited by
    \textcite{aguerayarcas2025intelligence} as holding that AI is not
    intelligent in the way we are. Read the two together.
  \item[Ewan Birney (EMBL-EBI)] Director of the European Bioinformatics
    Institute; a clear voice against hype from someone who lived through the
    genomics boom \parencite{birney2026genomics}.
  \item[Boris Cherny (Anthropic)] One of the creators of Claude Code, which they
    describe as ``a Unix utility rather than a traditional product'', made of
    the smallest pieces that are ``useful, understandable, and extensible''
    \parencite{liu2026claudecode} --- the argument of \Cref{ch:landscape}, from
    the people building the tool.
  \item[Christopher Manning (Stanford)] Foundations of language processing;
    co-author of the best empirical test of retrieval's limits
    \parencite{magesh2025legalrag}.
  \item[Lars Juhl Jensen (Copenhagen)] Biomedical text mining with dictionaries
    and practical tools --- still the right approach for many problems.
  \item[Sophia Ananiadou (Manchester)] Biomedical text mining, in the department
    where Turing worked.
  \item[Anne Thessen] Biodiversity informatics and ontologies: the data layer
    all of this depends on.
  \item[Martin Kleppmann (Cambridge)] Local-first software and distributed
    systems \parencite{kleppmann2019localfirst}.
  \item[Hadley Wickham (Posit)] On agents and AI in data analysis, from someone
    who has thought hard about tools for scientists.
\end{description}

\section{Sources}

\textbf{Primary.} arXiv (\code{cs.CL}, \code{cs.AI}, \code{cs.IR}) and bioRxiv.
Do not try to follow them; search them when you have a question, with alerts on
narrow terms only. Once a year, read the AI Index \parencite{aiindex2026}, the
best single overview. \textbf{Blogs and newsletters.} A few careful writers test
things by hand and report what fails: Simon Willison on local models
\parencite{willison2026qwen38}, Stephen Turner on AI in genomics and biosecurity
\parencite{turner2026slop}. \textbf{Lists and tutorials.} ``Awesome'' lists on
GitHub for RAG and agents, the mailing lists of the tools you use, and ELIXIR's
TeSS portal for training; check when a list was last updated, since a stale one
is worse than none. \textbf{Documentation.} Hugging Face for models, Ollama's
library for what runs locally, and your harness's own documentation, which
changes faster than any guide about it.

\textbf{Community.} Your domain community is more useful than the AI one. Surveys
of computational capacity in environmental omics find that what holds groups
back is training and keeping staff, not tools
\parencite{theodosopoulou2022greek} --- a corrective when talk turns to
hardware. TDWG for biodiversity standards, ELIXIR for life-science computing,
GBIF and OBIS community channels, and the World Biodiversity Forum, where
\textcite{pasche2026biodivpmc} and \textcite{krug2026ipbes} both appeared: the
people working on AI \emph{for biodiversity} are a small, findable group.

\section{A routine}

Weekly, thirty minutes: skim one aggregator and note things to look up. Monthly,
two hours: try one thing on a real task of your own. Quarterly, half a day:
review your tools, update pinned versions, test a backup restore. Yearly, a day:
read the AI Index, revisit your group's data and disclosure policy, reassess
hardware. About forty hours a year; more starts competing with doing ecology.

The fastest way to learn is a small real project: a RAG system over your own
reprints, a classifier for your camera-trap images, an extraction pipeline for
species descriptions, an agent-assisted clean-up of an analysis you have been
avoiding. You will find that PDF extraction is the hard part, that chunking
matters more than the model, and that the model is confidently wrong in ways you
can now recognise. None of that transfers from a blog post.

\section{What to tell colleagues}

\begin{quote}
These systems are useful for transforming text and writing code, where you
supply the content or can cheaply check the result. They are unreliable for
facts: they invent citations and agree with you when you are wrong. The main
change to daily work is agentic coding, which is closer to system administration
than to programming. The main thing to get right is data organisation and
provenance, which was worth getting right anyway. Run models locally where you
can, verify everything that matters, and disclose what you used.
\end{quote}

\section{A closing note}

The pressure to adopt is institutional, and it is not always aligned with good
work. It is legitimate to decide that an AI method does not suit a problem;
\textcite{mishra2026notready} makes that case for ecology seriously and well. It
is equally legitimate to use these tools heavily, with discipline. What is not
reasonable is using them without checking, or staying out of the argument while
others decide about your data, your infrastructure and your field.

Ayad Akhtar's play \emph{McNeal} \parencite{akhtar2024mcneal} follows a
celebrated novelist who feeds Sophocles, Shakespeare, Flaubert, Ibsen and his
late wife's journals into a chatbot and lets it help write his next novel.
Partway through he admits, ``my mind works better now if there's something
already on the page.'' Read by a scientist, the play is less about whether the
machine writes well than about what happens to the person who stops writing.
That is the risk to guard against, more than any single wrong answer.

The skills that last are the ones ecology already values: knowing where data
came from, doubting a result that is too convenient, and recording what you did
so that someone else can check it. Those transfer. Everything else is a lookup.
